固定1,024词表的实验显示，数学原子能够提高完整命令的表示，结构信号决定压缩与复合节点边界之间的工作点。Boundary AST在全领域新来源测试中取得0.3799的T/B，AST-v4取得73.33\%的复合节点对齐率和32.42 MiB/s的完整编码速度。四类划分将来源泛化、学科迁移与复杂度迁移分开测量，八项指标从压缩、词表使用、结构和运行效率描述各模型。下一阶段将采用相同模型架构、训练字节量和计算预算，从零训练接入不同分词器的模型，评价PDE／ODE求解、符号求导与积分、非方程类数学推理；同时比较训练吞吐、总训练计算量与任务准确率。
